\documentclass[11pt,a4paper]{article}

\usepackage[utf8]{inputenc}
\usepackage[T1]{fontenc}
\usepackage{booktabs}
\usepackage[margin=2.5cm]{geometry}
\usepackage{hyperref}
\usepackage{xcolor}
\usepackage{enumitem}

\hypersetup{
  colorlinks=true,
  linkcolor=blue!60!black,
  urlcolor=blue!60!black,
}

\title{Removing AI Watermarks with Off-the-Shelf Diffusion Models:\\
       A Practical Demonstration}
\author{Diffusion Watermarking Project}
\date{March 2026}

\begin{document}
\maketitle

%----------------------------------------------------------------------
\section{What Is a Diffusion Model?}

A diffusion model is a neural network that learns to generate images by
practising a very specific skill: starting from pure random noise and
gradually turning it into a coherent picture, one small step at a time.
Training works in reverse --- the model is shown real images with
increasing amounts of noise added, and learns to predict how to remove
the noise at each step. Once trained, you can hand it pure static and it
will denoise its way to a photorealistic image.

The key capability we exploit is called \textbf{image-to-image} (img2img)
regeneration. Instead of starting from pure noise, you start from an
existing image, add a controlled amount of noise (destroying some of the
original signal), and then let the model denoise it back into a clean
image. The result looks similar to the original --- same composition,
colours, and subject --- but every pixel has been \emph{regenerated} by
the diffusion model. The ``strength'' parameter controls how much noise
is added: low strength preserves more of the original; high strength
gives the model more creative freedom.

This is not a research prototype. Stable Diffusion is free, open-source
software that runs on a consumer laptop. The img2img pipeline is a
single function call in the standard library.

%----------------------------------------------------------------------
\section{The Attack}

Digital watermarks work by hiding a signal in the pixel values of an
image. The signal is imperceptible to the human eye but detectable by
the right algorithm. The watermark's security assumption is that an
attacker cannot remove this signal without visibly degrading the image.

The img2img attack breaks this assumption. When a watermarked image is
passed through the diffusion model:

\begin{enumerate}[leftmargin=2em]
  \item \textbf{Noise destroys the watermark.} The forward diffusion
    step adds Gaussian noise to the image, directly corrupting the
    learned signal pattern that the watermark encoder embedded.
  \item \textbf{Denoising reconstructs the image, not the watermark.}
    The model reconstructs a visually similar image from the noisy
    version, but the reconstruction is guided by the diffusion model's
    own learned prior --- not by the watermark encoder. The output
    carries no trace of the original watermark.
\end{enumerate}

No training or fine-tuning is required. The attacker does not need to
know anything about the watermark's architecture, key, or detection
method. A single forward pass through an unmodified, off-the-shelf
Stable Diffusion 1.5 checkpoint is sufficient.

%----------------------------------------------------------------------
\section{Results: VideoSeal (Meta)}

We first tested the attack against \textbf{VideoSeal}, Meta's
open-source neural watermark. VideoSeal embeds a 256-bit message into
images using a learned neural encoder/decoder --- it is the closest
publicly available analogue to Google's SynthID. We embedded watermarks
at an above-default strength (\texttt{scaling\_w}~$= 0.3$, 50\% above
the default) into 5 images from DiffusionDB, then attacked each at 3
strength levels with 2 guidance scales (6 combinations per image, 30
runs total).

\begin{table}[h]
\centering
\caption{VideoSeal attack results, averaged over 5 images. Bit accuracy
at chance level (50\%) means no watermark signal remains.}
\label{tab:videoseal}
\begin{tabular}{lcccc}
\toprule
\textbf{Attack Strength} & \textbf{SSIM} & \textbf{PSNR (dB)} &
\textbf{Bit Accuracy} & \textbf{Watermark Detected} \\
\midrule
\textit{Before attack} & \textit{0.991} & \textit{41.6} &
\textit{100.0\%} & \textit{5/5} \\
\midrule
0.3 (mild)     & 0.540 & 20.4 & 53.1\% & 0/5 \\
0.5 (moderate) & 0.433 & 18.2 & 49.9\% & 0/5 \\
0.7 (strong)   & 0.337 & 16.0 & 49.6\% & 0/5 \\
\midrule
\textit{Chance level} & --- & --- & \textit{50.0\%} & --- \\
\bottomrule
\end{tabular}
\end{table}

\textbf{The watermark was removed in 100\% of trials} (30/30), at every
attack strength tested. Even the mildest setting (strength 0.3) reduced
bit accuracy from 100\% to 53\% --- statistically indistinguishable from
random guessing. The $p$-value collapsed from $10^{-78}$ (before attack)
to $\sim\!0.1$ (after), a drop of 77 orders of magnitude.

The attacked images are visibly altered (SSIM drops to 0.34--0.54) but
remain recognisable. At strength 0.3, image quality is roughly
comparable to moderate JPEG compression.

%----------------------------------------------------------------------
\section{Pilot Test: SynthID (Google)}

Having established the attack on VideoSeal, we turned to the real
target: \textbf{SynthID}, Google DeepMind's 136-bit watermark deployed
across all Gemini-generated images. SynthID's detector is closed-source,
but Google provides a public detection portal that returns a binary
yes/no verdict.

For the pilot, we generated a single image via the Gemini API (``old man
in a boat'', $1408 \times 768$). Because Stable Diffusion operates at
$512 \times 512$, we used a \textbf{tiled attack}: split the image into
overlapping patches, attack each tile independently, then reassemble.
We tested three strengths and uploaded each result to Google's detection
portal.

\begin{table}[h]
\centering
\caption{Pilot SynthID attack ($n = 1$). Detection via Google's
SynthID portal.}
\label{tab:synthid-pilot}
\begin{tabular}{lccc}
\toprule
\textbf{Attack Strength} & \textbf{SSIM} & \textbf{PSNR (dB)} &
\textbf{SynthID Detected} \\
\midrule
\textit{Original image} & \textit{1.000} & $\infty$ & \textit{Yes} \\
\midrule
0.2 & 0.684 & 22.1 & \textbf{No} \\
0.3 & 0.617 & 20.7 & \textbf{No} \\
0.5 & 0.472 & 18.3 & \textbf{No} \\
\bottomrule
\end{tabular}
\end{table}

All three attack strengths removed the SynthID watermark according to
Google's own detector. This is only $n = 1$, but it demonstrates that
the attack concept transfers from VideoSeal to SynthID without any
modification.

%----------------------------------------------------------------------
\section{Scaling Up: 50-Image SynthID Corpus}

To move beyond a single data point, we generated a \textbf{50-image
corpus} via the Gemini API using diverse prompts (landscapes, portraits,
abstract art, food, sci-fi, etc.; all $1024 \times 1024$ PNG). Each
image is attacked at six fine-grained strengths using the same tiled
img2img pipeline.

Since we cannot submit 300 images to Google's portal manually, we use
\textbf{reverse-SynthID}, an open-source detector that analyses
Fourier-domain carrier frequencies characteristic of SynthID. While less
authoritative than Google's detector, it provides continuous metrics
(correlation strength, phase coherence) rather than just a binary
verdict, letting us track exactly how the watermark signal degrades.

\begin{table}[h]
\centering
\caption{50-image SynthID corpus sweep. Results pending --- sweep in
progress at time of writing.}
\label{tab:synthid-sweep}
\begin{tabular}{lcccc}
\toprule
\textbf{Attack Strength} & \textbf{Mean SSIM} & \textbf{Mean PSNR} &
\textbf{Mean Correlation} & \textbf{Detection Rate} \\
\midrule
0.05 & --- & --- & --- & ---/50 \\
0.10 & --- & --- & --- & ---/50 \\
0.15 & --- & --- & --- & ---/50 \\
0.20 & --- & --- & --- & ---/50 \\
0.25 & --- & --- & --- & ---/50 \\
0.30 & --- & --- & --- & ---/50 \\
\midrule
\textit{Original} & \textit{1.000} & $\infty$ & \textit{TBD} &
\textit{50/50} \\
\bottomrule
\end{tabular}
\end{table}

Preliminary results on 2 images show that correlation drops by
$\sim$\!3$\times$ between the original and strength 0.3, confirming the
watermark signal is being degraded. The full sweep ($50 \times 6 = 300$
attacks) will reveal whether SynthID has a sharp breaking point or a
gradual decline.

%----------------------------------------------------------------------
\section{What This Means}

The core finding is simple: \textbf{an off-the-shelf diffusion model,
used as-is, removes state-of-the-art neural watermarks.}

\begin{itemize}[leftmargin=2em]
  \item \textbf{No expertise required.} The attack is a single function
    call in a free, open-source library. Anyone who can install Python
    can run it.
  \item \textbf{No knowledge of the watermark needed.} The attacker
    does not need to know the watermark's architecture, key, embedding
    strength, or detection algorithm. The diffusion model is not
    ``trained against'' the watermark --- it simply destroys the signal
    as a side effect of its normal operation.
  \item \textbf{Transfers across watermark systems.} The same attack
    (same model, same code, same parameters) removed both VideoSeal
    (Meta, 256-bit) and SynthID (Google, 136-bit) watermarks, despite
    having no access to either system's internals.
  \item \textbf{Quality--robustness tradeoff favours the attacker.}
    Even the mildest attack setting (strength 0.3) completely removes
    the watermark while preserving recognisable image content. The
    attacker can choose how much quality to sacrifice.
\end{itemize}

This does not mean watermarking is useless --- it raises the bar for
casual redistribution and provides a signal for platforms that control
the upload pipeline. But it does mean that any motivated individual can
strip the watermark in seconds, which limits the technology's value as
a forensic provenance tool against adversarial actors.

The watermark cannot survive regeneration because it is embedded in
pixel-level statistics, and the diffusion model regenerates those
statistics from scratch. Any watermark that operates in the same signal
domain as the diffusion model's output is fundamentally vulnerable to
this class of attack.

\end{document}
